\section{Pushing paths and local connectedness}\label{sec:path-pushing}

We work in the setting of Proposition~\ref{block:setting}, with compact
metric space $Y=H\sqcup L$ and graph $K$ on $I$, and use the annulus system
of Lemma~\ref{ann:system}.  Our aim is to turn paths near a point of $L$
into connected subsets of $L$.  To do this we push the paths arbitrarily
far from a fixed closed set and then take limits of their vertex sets.
The graph can have infinite valence and infinite vertex stabilizers;
none of the arguments below requires finite generation of $H$.

At a fixed depth, nearby boundary points determine components of the
remaining graph. We will prove that, beyond one common threshold,
increasing the depth cannot separate points with the same component
label. Iterating this fact produces a single boundary neighborhood whose
points retain the same component label at every greater depth. Limits
of the resulting graph paths will then lie entirely in the boundary.

\subsection{A uniform depth comparison}

Fix a closed subset $A\subset Y$ with $1\in A$ and $L\setminus A\ne\varnothing$.
The annular depth $D(F,v)$ is the maximum length of a chain from $F$ to $v$
when it is finite, as in Lemma~\ref{ann:depth}.  Write
\[
 d_0=\operatorname{dist}(1,L)>0,
 \qquad
 N(a,b)=\#\{\sigma:
       \operatorname{diam}\sigma^-\ge a,
       \operatorname{diam}\sigma^+\ge b\}\quad(a,b>0).
\]
The number $N(a,b)$ is finite by Lemma~\ref{ann:system}.  A side containing
$1$ has diameter at least $d_0$, because every side meets $L$.

\begin{lemma}\label{push:shrinking}
Let $g_n\in H$ satisfy $D(A,g_n)\to\infty$, and set $F_n=g_n^{-1}A$.
After passage to a subsequence there is $q\in L$ such that $F_n$ converges
in the Hausdorff metric to $\{q\}$.
\end{lemma}
\begin{proof}
By equivariance, $D(F_n,1)=D(A,g_n)$.  In any chain from $F_n$ to $1$,
every positive side contains $1$.  If the chain has more than
$N(\epsilon,d_0)$ members, one of its negative sides has diameter less
than $\epsilon$.  This side contains $F_n$.  Hence
$\operatorname{diam}F_n\to0$.

Lemma~\ref{ann:depth} gives $D(A,g)<\infty$ for each fixed $g\in H$.
Thus $g_n$ escapes every finite subset of $H$, as does $g_n^{-1}$.
Since $1\in A$, we have $g_n^{-1}\in F_n$.  Compactness of $Y$ and
isolatedness of its group vertices give a subsequence
$g_n^{-1}\to q\in L$.  The diameter conclusion now gives the assertion.
\end{proof}

\begin{lemma}\label{push:depth-comparison}
Suppose $F_n=g_n^{-1}A$ converges to $\{q\}$ as in
Lemma~\ref{push:shrinking}.  For each $r>0$, set
\[
 C_r=N(r/2,d_0)+N(d_0,r)+2.
\]
For all sufficiently large $n$, uniformly over every $b\in H$ with
$\rho(b,q)\ge3r$, one has
\begin{equation}\label{eq:push-depth-comparison}
 D(F_n,b)\ge D(F_n,1)+D(\{1\},b)-C_r.
\end{equation}
Consequently, if $b_n\in H$ tends to $z\in L\setminus\{q\}$, then
\[
 D(F_n,b_n)\ge D(F_n,1)+D(\{1\},b_n)-C
 \quad\hbox{and}\quad D(\{1\},b_n)\longrightarrow\infty
\]
for a constant $C$ independent of $n$.
\end{lemma}
\begin{proof}
Choose $n$ large enough that
$F_n\subset B(q,r/2)$ and
$\ell=D(F_n,1)>N(r/2,d_0)+1$.  Take a chain
$\sigma_1<\cdots<\sigma_\ell$ of maximum length from $F_n$ to $1$,
and put $U=B(q,r)$.
If $\sigma_j^-\not\subset U$, its negative side contains a point of
$F_n$ and a point outside $U$, so has diameter at least $r/2$.
Its positive side has diameter at least $d_0$.
There are therefore at most $N_-=N(r/2,d_0)$ such annuli.
Because negative sides increase along a chain, these form a terminal
segment.  Remove that segment and one further annulus.
The retained initial chain has at least $\ell-N_--1$ members.
The negative side of the next annulus lies in $U$, so the last retained
positive side contains $Y\setminus U$.
In particular, this retained chain already ends on the side containing
any $b$ under consideration.

Let $t=D(\{1\},b)$, and take a maximum chain
$\tau_1<\cdots<\tau_t$ from $1$ to $b$ when $t>0$.
Put $V=Y\setminus\overline{B(q,2r)}$.
A positive side not contained in $V$ contains both $b$ and a point at
distance at most $2r$ from $q$, so has diameter at least $r$.
Every negative side contains $1$, so has diameter at least $d_0$.
Thus at most $N_+=N(d_0,r)$ positive sides fail to lie in $V$;
these form an initial segment.
If $t>N_++1$, remove this initial segment and one further annulus.
The retained final chain has at least $t-N_+-1$ members, and its first
negative side contains $Y\setminus V$.

Since $U\cap V=\varnothing$, the last positive side of the first retained
chain and the first negative side of the second cover $Y$.
The defining order relation therefore concatenates the two chains.
Their total length is at least $\ell+t-N_--N_+-2$.
If $t\le N_++1$, the first retained chain alone has at least this length.
This proves \eqref{eq:push-depth-comparison}, including a constant and
threshold independent of $b$.

For the final assertion choose $r$ with $4r<\rho(q,z)$.
Eventually $\rho(b_n,q)\ge3r$, so the inequality applies.
For any prescribed integer $N$, Lemma~\ref{ann:system} supplies a chain
of length $N$ from $\{1\}$ to a neighborhood of $z$.
Eventually this neighborhood contains $b_n$, proving
$D(\{1\},b_n)\to\infty$.
\end{proof}

Figure~\ref{fig:annular-chains} illustrates this pruning and
concatenation for a sequence $b_n\to z\ne q$.
\begin{figure}[!htbp]
\centering
\begin{tikzpicture}[
  x=1cm,y=1cm,
  every node/.style={font=\small},
  retained/.style={draw,rounded corners=2pt,fill=black!4,
    minimum height=0.68cm,inner xsep=8pt},
  omitted/.style={draw=black!70,dashed,rounded corners=2pt,
    text=black!85,minimum height=0.68cm,inner xsep=7pt},
  link/.style={-{Stealth[length=1.7mm]},semithick},
  discarded/.style={link,dashed,draw=black!70}
]
\node (f) at (0,0) {$F_n$};
\node[retained] (sig) at (3.1,0)
  {$\sigma_1<\cdots<\sigma_i$};
\node[omitted] (tail) at (7.25,0) {discard $\le N_-+1$};
\node (onea) at (10.1,0) {$1$};
\draw[link] (f) -- (sig);
\draw[discarded] (sig) -- (tail);
\draw[discarded] (tail) -- (onea);

\node (oneb) at (0,-1.2) {$1$};
\node[omitted] (head) at (2.85,-1.2) {discard $\le N_++1$};
\node[retained] (tau) at (7,-1.2)
  {$\tau_j<\cdots<\tau_t$};
\node (b) at (10.1,-1.2) {$b_n$};
\draw[discarded] (oneb) -- (head);
\draw[discarded] (head) -- (tau);
\draw[link] (tau) -- (b);

\node[align=center] at (5.05,-2.48) {
  $Y\setminus U\subset\sigma_i^+\,,\qquad
   U\cap V=\varnothing\,,\qquad
   Y\setminus V\subset\tau_j^-$\\[5pt]
  $\sigma_i^+\cup\tau_j^-=Y\quad\Longrightarrow\quad
   \sigma_i<\tau_j$
};
\node[retained,inner xsep=12pt] (joined) at (5.05,-3.82)
  {$\sigma_1<\cdots<\sigma_i<\tau_j<\cdots<\tau_t$};
\node (fj) at (0,-3.82) {$F_n$};
\node (bj) at (10.1,-3.82) {$b_n$};
\draw[link] (fj) -- (joined);
\draw[link] (joined) -- (bj);
\end{tikzpicture}
\caption{Concatenating annular chains after bounded pruning. Choose
fixed disjoint neighborhoods $U$ of $q$ and $V$ of $z$, where
$F_n$ shrinks to $q$ and $b_n\to z\ne q$.
The terminal positive side of the first retained chain contains
$Y\setminus U$; the initial negative side of the second contains
$Y\setminus V$. Their union is $Y$, which is exactly the annulus-order
criterion for joining the chains. The diagram records order and
containment only; the sides need not be connected.}
\label{fig:annular-chains}
\end{figure}


\subsection{Components of the deeper vertex sets}

Choose one representative $p_i$ for each of the finitely many
$H$-orbits in $I$.  When $I=H$, use only $p_1=1$.
For a vertex $u$ in the orbit of $p_i$, define its set of representatives
in the group by
\[
 C(u)=\{g\in H:g p_i=u\}.
\]
For integers $m\ge0$ define
\begin{equation}\label{eq:push-deep-vertices}
 I_m=\left\{u\in I\setminus A:
           \min_{g\in C(u)}D(A,g)\ge m\right\}.
\end{equation}
Each depth in this minimum is a finite nonnegative integer, so the
minimum exists.  The condition concerns \emph{every} representative of
$u$.  This will allow us to move through an infinite stabilizer without
losing the depth bound.

\begin{lemma}\label{push:component-labels}
For each fixed $m$, the deleted vertices $I\setminus I_m$ have no
accumulation by distinct vertices at any point of $L\setminus A$.
Consequently Lemma~\ref{ann:component} defines a locally constant map
\[
 Q_m:L\setminus A\longrightarrow
      \{\text{components of }K[I_m]\},
\]
whose value at $x$ contains all vertices of $I_m$ sufficiently near $x$.
Every $x\in L\setminus A$ is approached by distinct vertices with this
label.  Moreover, for each finite subset $E\subset I$ one has
$E\cap I_m=\varnothing$ for all sufficiently large $m$.
\end{lemma}
\begin{proof}
For $m=0$ the first assertion follows from $I_0=I\setminus A$.
Suppose $m\ge1$.
Suppose that distinct $u_n\in I\setminus I_m$ tend to
$x\in L\setminus A$.
Eventually $u_n\notin A$, so choose $a_n\in C(u_n)$ with $D(A,a_n)<m$.
Since there are only finitely many orbit representatives and the $u_n$
are distinct, $a_n$ escapes every finite subset of $H$.
The comparison $\rho(gp_i,g)\to0$ in Proposition~\ref{block:setting},
applied to the finite list of representatives, gives $a_n\to x$.
A chain from $A$ to a neighborhood of $x$ of length $m$, supplied by
Lemma~\ref{ann:system}, contradicts $D(A,a_n)<m$ for large $n$.
The assertions about $Q_m$ and approximation follow from
Lemma~\ref{ann:component}.

Finally, for each $u\in E\setminus A$ choose a single $a_u\in C(u)$.
The depth $D(A,a_u)$ is finite.  Taking $m$ larger than all these finitely
many depths excludes every vertex of $E$ from $I_m$.
\end{proof}

Our goal is now to find one depth threshold beyond which increasing
$m$ to $m+1$ never separates boundary points with the same component
label. The next three lemmas establish the uniformity needed for
Lemma~\ref{push:refinement}.

The next lemma is the main use of the uniform depth estimate.
It says that translating from a vertex of large depth leaves only
finitely many deleted vertices in any fixed region away from one
boundary point.  The exceptional finite set is independent of the
index of the sequence.

\begin{lemma}\label{push:translated-deletions}
Let $m_n\to\infty$ and $D(A,g_n)\ge m_n$.
Pass to a subsequence for which $F_n=g_n^{-1}A$ tends to $\{q\}$,
with $q\in L$.
For every compact $E\subset Y\setminus\{q\}$ there are a finite
set $S_E\subset I$ and an index $n_E$ such that
\[
 E\cap g_n^{-1}(I\setminus I_{m_n+1})\subset S_E
 \qquad(n\ge n_E).
\]
\end{lemma}
\begin{proof}
If no such finite set and index exist, choose indices tending to infinity
and pairwise distinct vertices; relabel them so that
\[
 w_n\in E\cap g_n^{-1}(I\setminus I_{m_n+1}).
\]
Compactness gives a further subsequence $w_n\to z\in E$.
The limit lies in $L$: this is automatic when $I\subset L$, and follows
from isolatedness of group vertices when $I=H$.
Since $F_n$ tends to $q\notin E$, eventually $w_n\notin F_n$.
The deletion of $g_nw_n$ therefore gives a representative
$g_nb_n\in C(g_nw_n)$ such that
\[
 D(F_n,b_n)=D(A,g_nb_n)<m_n+1.
\]
After restricting to one orbit representative we have $w_n=b_np_i$.
Distinctness of the $w_n$ forces $b_n$ to escape finite subsets of $H$,
and the comparison with group vertices gives $b_n\to z\ne q$.
Lemma~\ref{push:depth-comparison} now yields
\[
 D(F_n,b_n)\ge m_n+D(\{1\},b_n)-C>m_n+1
\]
for large $n$, a contradiction.
\end{proof}

\begin{lemma}\label{push:eventual-labels}
For the sequences in Lemma~\ref{push:translated-deletions}, define
\[
 \ell_n(z)=Q_{m_n+1}(g_nz)
\]
whenever $z\in L$ and $g_nz\notin A$.
For each $z\in L\setminus\{q\}$ there are a relative neighborhood
$V_z\subset L\setminus\{q\}$ and an index $n_z$ such that all values
$\ell_n(w)$ with $w\in V_z$ and $n\ge n_z$ are defined and
\[
 \ell_n(w)=\ell_n(z)\qquad(w\in V_z,\ n\ge n_z).
\]
For any two points $z,z'\in L\setminus\{q\}$ one consequently has
$\ell_n(z)=\ell_n(z')$ for all sufficiently large $n$.
The last threshold may depend on the two points.
\end{lemma}
\begin{proof}
Choose an open neighborhood $O$ of $z$ in $Y$ whose closure misses $q$.
Apply Lemma~\ref{push:translated-deletions} to the compact set
$E=\overline O$. There are a finite set
$S\subset I$ and an index $n_0$ such that
\[
 T=(I\cap O)\setminus S\quad\hbox{satisfies}\quad
 g_nT\subset I_{m_n+1}\qquad(n\ge n_0).
\]
The deleted vertices $I\setminus T$ have no distinct accumulation at any
point of $O\cap L$.  Lemma~\ref{ann:component} therefore gives a
neighborhood $V_z\subset O\cap L$ on which the component of $K[T]$
approached by vertices is constant; denote it by $C$.
Shrink $V_z$ if needed.  Since $F_n\to\{q\}$, all $g_nw$ for
$w\in V_z$ lie outside $A$ once $n$ is large.

Fix such an $n$.  Every $w\in V_z$ is approached by vertices of $C$.
Their images approach $g_nw$ and lie in the connected subgraph $g_nC$
of $K[I_{m_n+1}]$.
The defining property of $Q_{m_n+1}$ therefore identifies
$\ell_n(w)$ with the component containing $g_nC$, for every $w\in V_z$.
This gives the asserted neighborhood and threshold, simultaneously for
all its points.

On $L\setminus\{q\}$, declare $z\sim z'$ if their labels agree for all
sufficiently large $n$.  Eventual equality is an equivalence relation,
even though the label set can change with $n$.
The first part shows that every equivalence class is open.
Since $q$ is not a cut point, $L\setminus\{q\}$ is connected, so this
partition has one class.
\end{proof}

\subsection{A common threshold for all graph paths}

So far the conclusions concerned a chosen sequence of deep group
vertices.  We now use three fixed boundary points to obtain one
threshold that works for every graph path.  The number three ensures
that two of these points always avoid the single exceptional point in
Lemma~\ref{push:eventual-labels}.

Choose a set $Z\subset L$ of three distinct points; it exists because
$L$ is a nondegenerate continuum.

\begin{lemma}\label{push:majority}
For every finite set $S\subset H$, there is an integer $M=M(A,Z,S)\ge1$ with the
following properties for all integers $m\ge M$.
\begin{enumerate}
\item If $D(A,g)\ge m$, at least two points of $gZ$ lie in
$L\setminus A$ and have a common $Q_{m+1}$ label.  Denote this unique
majority label by $J_m(g)$.
\item If $s\in S$ and $D(A,g),D(A,gs)\ge m$, then
$J_m(g)=J_m(gs)$.
\end{enumerate}
\end{lemma}
\begin{proof}
If the first assertion fails at arbitrarily large depths, choose
$m_n\to\infty$ and $g_n$ violating it with $D(A,g_n)\ge m_n$.
Pass to a subsequence with $g_n^{-1}A\to\{q\}$.
Two fixed points of $Z$ avoid $q$.
By Lemma~\ref{push:eventual-labels}, their images under $g_n$ eventually
lie outside $A$ and have the same $Q_{m_n+1}$ label, a contradiction.
This also proves uniqueness of the majority label, since two different
labels cannot each occur at two points of a three-point set.

After choosing a threshold for the first assertion, suppose that the
second fails at arbitrarily large $m$.
Choose a violating sequence and, since $S$ is finite, restrict to one
fixed $s\in S$.
Again pass to a subsequence with $g_n^{-1}A\to\{q\}$.
Choose two points of $Z\setminus\{q\}$ and two points of
$sZ\setminus\{q\}$.
Lemma~\ref{push:eventual-labels} makes the labels of all four images
under $g_n$ equal for all sufficiently large $n$.
They determine the majority labels of both $g_nZ$ and $g_nsZ$,
contradicting the chosen inequality.
\end{proof}

We choose the finite set $S$ once and for all as follows.
There are finitely many orbits of ordered edges: allowing both
orientations at most doubles the number of edge orbits.
For each representative $(u_e,v_e)$ choose $a_e,b_e\in H$ and indices
$i(e),j(e)$ such that
\[
 u_e=a_ep_{i(e)},\qquad v_e=b_ep_{j(e)},
\]
and include $a_e^{-1}b_e$ in $S$.
Thus the endpoints of every edge have representatives $g$ and $gs$
for some $s\in S$.
When $I\subset L$, also include, for each $p_i$, a stabilizing element
$h_i\in H$ supplied by Proposition~\ref{block:setting}, with
\[
 h_i p_i=p_i,\qquad h_i^j z\longrightarrow p_i
 \quad(z\in L,\ j\longrightarrow+\infty).
\]
Fix the threshold $M$ of Lemma~\ref{push:majority} for this list.

\begin{lemma}\label{push:refinement}
For every $m\ge M$ and all $x,y\in L\setminus A$,
\begin{equation}\label{eq:push-refinement}
 Q_m(x)=Q_m(y)\quad\Longrightarrow\quad
 Q_{m+1}(x)=Q_{m+1}(y).
\end{equation}
In particular, the threshold $M$ is independent of $x$ and $y$.
\end{lemma}
\begin{proof}
First suppose $I\subset L$.
We claim that for every $v\in I_m$ and every $g\in C(v)$,
\begin{equation}\label{eq:push-support-label}
 J_m(g)=Q_{m+1}(v).
\end{equation}
Let $v=gp_i$.
For all $j\ge0$, the element $gh_i^j$ also belongs to $C(v)$.
The minimum in \eqref{eq:push-deep-vertices} gives
$D(A,gh_i^j)\ge m$ for every $j$.
The move $h_i$ belongs to $S$, so Lemma~\ref{push:majority} gives
\[
 J_m(g)=J_m(gh_i)=J_m(gh_i^2)=\cdots.
\]
Here $g$ and $m$ are fixed.
As $j\to\infty$, the three anchors $gh_i^jZ$ tend to $v$.
Since $v\in L\setminus A$ and $Q_{m+1}$ is locally constant there,
all three anchors eventually have label $Q_{m+1}(v)$.
This proves \eqref{eq:push-support-label}.
The label at $v$ is defined even when $v\notin I_{m+1}$;
it records the component approached by allowed vertices near $v$.

If $v,w\in I_m$ are adjacent, choose their representatives $g,gs$
using the finite edge list.
Both depths are at least $m$ by the minimum convention.
Lemma~\ref{push:majority} and \eqref{eq:push-support-label} give
$Q_{m+1}(v)=Q_{m+1}(w)$.
Thus $Q_{m+1}$ is constant along every path in $K[I_m]$.
If $Q_m(x)=Q_m(y)$, choose vertices in this component sufficiently near
$x$ and $y$ that their $Q_{m+1}$ labels equal those of $x$ and $y$.
Such vertices exist by Lemma~\ref{push:component-labels} and local
constancy of $Q_{m+1}$.
A path in their common component proves
\eqref{eq:push-refinement}.

Now suppose $I=H$.
There is a single representative $p_1=1$, and $C(v)=\{v\}$.
The finite edge moves and Lemma~\ref{push:majority} imply that $J_m$
is constant along paths of $K[I_m]$.
For a fixed $x\in L\setminus A$ and fixed $m\ge M$, we claim that
\[
 J_m(g)=Q_{m+1}(x)
 \quad\hbox{for all group vertices $g$ sufficiently near $x$}.
\]
These vertices belong to $I_m$ when sufficiently near $x$, by
Lemma~\ref{push:component-labels}.
If the claim fails, take a sequence $g_n\to x$ with the wrong label,
and pass to a subsequence $g_n^{-1}\to b\in L$.
The convergence property in Proposition~\ref{block:setting} sends at
least two fixed points of $Z\setminus\{b\}$ to $x$ under $g_n$.
Local constancy of $Q_{m+1}$ near $x$ gives the asserted majority label,
a contradiction.
Approximating $x$ and $y$ inside their common $Q_m$ component and using
constancy of $J_m$ along a connecting path now proves
\eqref{eq:push-refinement} in this case as well.
\end{proof}

\subsection{Limits of the pushed paths}

\begin{proposition}\label{push:local-connected}
In the setting of Proposition~\ref{block:setting}, the continuum $L$
is locally connected.
\end{proposition}
\begin{proof}
Fix $x\in L$ and a relatively open neighborhood $U\subset L$ of $x$.
Choose an open neighborhood $O$ of $x$ in $Y$ with $1\notin O$ and
$\overline O\cap L\subset U$, and apply the preceding construction to
$A=Y\setminus O$.
Let $M$ be the threshold in Lemma~\ref{push:refinement}.
By local constancy of $Q_M$, there is a relative neighborhood
$V\subset L\setminus A$ of $x$ such that $Q_M(y)=Q_M(x)$ for all
$y\in V$.
Induction using Lemma~\ref{push:refinement} gives
\[
 Q_m(y)=Q_m(x)\qquad(y\in V,\ m\ge M).
\]

Fix $y\in V$.
For each $m\ge M$ choose vertices $a_m,b_m\in I_m$ in this common
component with $\rho(a_m,x)<1/m$ and $\rho(b_m,y)<1/m$,
and choose a finite path $P_m$ in $K[I_m]$ joining them.
Every vertex of $P_m$ lies in $O$.
We next show that the maximum distance between consecutive vertices of
these paths tends to zero.

Use the finite list of edge representatives $(u_e,v_e)$.
The comparison with group vertices gives
\[
 \rho(gu_e,gv_e)
 \le\rho(gu_e,g)+\rho(g,gv_e)\longrightarrow0
 \qquad(g\longrightarrow\infty\hbox{ in }H),
\]
uniformly over this finite list.
For each $\epsilon>0$, choose a finite $F\subset H$ outside which every
one of these edge distances is less than $\epsilon$.
There are only finitely many endpoints of the edges
$(gu_e,gv_e)$ with $g\in F$.
By Lemma~\ref{push:component-labels}, none of these endpoints belongs to
$I_m$ once $m$ is large.
Consequently every edge in $K[I_m]$ then has endpoint distance less than
$\epsilon$.  This proves the mesh assertion without any local-finiteness
assumption on $K$.

Regard $P_m$ as its finite vertex set.
The space of nonempty closed subsets of the compact metric space $Y$
is compact in the Hausdorff metric, so pass to a subsequence with
$P_m\to C$.
The limit contains $x,y$ and lies in $\overline O$.
It also lies in $L$: this is immediate if $I\subset L$; if $I=H$,
Lemma~\ref{push:component-labels} excludes each fixed group vertex from
$P_m$ for large $m$, and isolatedness of that vertex excludes it from
the Hausdorff limit.

The limit $C$ is connected.
Otherwise write it as a union of disjoint nonempty compact sets $C_0,C_1$,
whose distance is some $\delta>0$.
For large $m$, the set $P_m$ lies in the union of their disjoint
$\delta/3$-neighborhoods and meets both.
A path with these vertices must have two consecutive vertices in
different neighborhoods, at distance at least $\delta/3$.
This contradicts the mesh assertion.
Thus $C$ is a connected subset of $\overline O\cap L\subset U$
containing $x$ and $y$.

Every $y\in V$ therefore belongs to the component of $x$ in $U$.
Repeating the argument at every point of every open subset of $L$ shows
that its components are open.  This is local connectedness of $L$.
\end{proof}
